\documentclass[10pt,a4paper]{amsart}
\usepackage[margin=19mm]{geometry}
\usepackage{amsmath,amssymb,mathtools}
\usepackage[scr=rsfs,cal=euler]{mathalfa}
\usepackage{xfrac}
\usepackage[unicode,hidelinks,bookmarks=false]{hyperref}
\newcommand{\ie}{i.e.\ }
\newcommand{\cf}{cf.\ }
\newcommand{\resp}{respectively\ }
\newcommand{\Subsection}{\subsection}
\providecommand{\nobreakdash}{\nobreak\hbox{-}}
\setlength{\emergencystretch}{2em}
\multlinegap0pt
\usepackage{bm}
\usepackage{bbm}
\usepackage{dsfont}
\usepackage{xspace}
\usepackage{cancel}
\usepackage{mathtools}
\usepackage{amsthm}
\makeatletter
\@ifundefined{theorem}{\newtheorem{theorem}{Theorem}}{}
\@ifundefined{lemma}{\newtheorem{lemma}[theorem]{Lemma}}{}
\makeatother
\title[Sums of the floor function related to class numbers of imagi]{Sums of the floor function related to class numbers of imaginary quadratic fields}
\author{Marc Chamberland, Karl Dilcher}
\date{Source: arXiv:2510.04387v1 (CC BY 4.0)}
\begin{document}
\maketitle

\noindent\emph{Source and licence.} Sums of the floor function related to class numbers of imaginary quadratic fields. Version arXiv:2510.04387v1 (CC BY 4.0), distributed under the Creative Commons Attribution 4.0 International licence, as recorded in the arXiv licence metadata for this exact version and shown on the abstract page https://arxiv.org/abs/2510.04387v1. Full rights notice: licenses/arxiv-2510.04387-CC-BY-4.0.txt.\par\smallskip

\noindent\textit{Editorial scope.} Counted unit: the Lemma whose source label is \texttt{lem:6.2} in the article (source0.tex lines 1056--1081, source label \texttt{lem:6.2}), delivered together with the article's own complete original proof of it (source0.tex lines 1196--1342, one \texttt{proof} environment of 147 lines). No mathematics is written, repaired or re-derived: statement and proof are the article's own text line for line. The proof is detached from the statement in the source: the article states the result at lines 1056--1081 and prints its proof at lines 1196--1342, and the proof environment's own header names the statement, so the association is the article's own. Required context retained uncounted: 2.24 (lines 433--436); 2.25 (lines 452--454); 3.2 (lines 594--596); 3.7 (lines 692--694); 4.11 (lines 902--905); lem:6.3 (lines 1130--1151); 6.12 (lines 1132--1135). Every definition, display and label of those lines is retained unchanged. Omitted from this package and not redistributed: 1--432, 437--451, 455--593, 597--691, 695--901, 906--1055, 1082--1129, 1136--1195, 1343--1544. Nothing inside a retained excerpt is omitted or altered: every retained body is delivered exactly as the article's lines are written, so no omission receipt of any kind accompanies this package. Citations: the retained excerpts carry no citation command at all, so no bibliography entry of the article is transcribed and no reference is redistributed. Reference adaptation: none was needed and none was made; every reference inside the delivered unit resolves inside it, so no unresolved reference or citation is produced. Printed slips kept literal: no printed word, symbol, hypothesis or number of the counted statement or of its proof was altered. Layout and class: the article's own class is \texttt{amsart} with packages including amssymb, graphicx; the wrapper is the lane's pinned \texttt{amsart} editorial wrapper at 10pt on A4, which restates the theorem-like environments the retained text needs and adds the article's own macro definitions that the retained text needs (none), with extra wrapper packages (bm, bbm, dsfont, xspace, cancel, mathtools). The delivered numbering is the unit's own. Assets: no retained span contains a figure environment, a graphics-inclusion command, a PDF-page inclusion, a table or a TikZ picture; no image file is copied into this package, the package declares no asset dependency and no image file is redistributed with it. Licence: version arXiv:2510.04387v1 of this article is distributed under the Creative Commons Attribution 4.0 International licence, as recorded in the arXiv licence metadata for this exact version and shown on the abstract page https://arxiv.org/abs/2510.04387v1. A full rights notice is delivered at licenses/arxiv-2510.04387-CC-BY-4.0.txt. Characters: every retained excerpt is pure ASCII, so the delivered engine, XeLaTeX with its default Latin Modern text font, reports no missing character.
\begin{equation}\label{2.24}
\sum_{k=1}^{\frac{p-1}{2}}{\rm Rem}(k^2\div p) = \frac{p(p-1)}{4}
+\frac{1}{2}\sum_{j=1}^{p-1}j\left(\frac{j}{p}\right).
\end{equation}

\begin{equation}\label{2.25}
\sum_{j=1}^{p-1}j\left(\frac{j}{p}\right)=0\qquad (p\equiv 1\pmod{4}).
\end{equation}

\begin{equation}\label{3.2}
h(K)=\frac{-w(d)}{2|d|}\sum_{j=1}^{|d|-1}j\left(\frac{d}{j}\right),
\end{equation}

\begin{equation}\label{3.7}
h(-p)=-\frac{1}{p}\sum_{j=1}^{p-1}j\left(\frac{j}{p}\right),
\end{equation}

\begin{equation}\label{4.11}
\sum_{k=1}^{p-1}{\rm Rem}(k^2\div p) = \frac{p(p-1)}{2}
-p\cdot h(-p)\qquad (p\equiv 3\pmod{4}),
\end{equation}

\begin{lemma}\label{lem:6.3}
Let $p$ be a prime. $(a)$ For any integers $k\geq 1$ and $\alpha\geq 1$ we have
\begin{equation}\label{6.12}
{\rm Rem}\left((kp)^2\div p^{\alpha+2}\right)
=p^2\cdot{\rm Rem}\left(k^2\div p^{\alpha}\right).
\end{equation}
$(b)$ For fixed integers $r$, $1\leq r\leq p-1$, and $\alpha\geq 0$, the 
following remainders are distinct:
\begin{equation}\label{6.13}
{\rm Rem}\left((jp+r)^2\div p^{\alpha+2}\right),\qquad
0\leq j\leq p^{\alpha+1}-1.
\end{equation}
$(c)$ For $\alpha\geq 2$ we have
\begin{equation}\label{6.14}
{\rm Rem}((2^{\alpha-1}-k)^2\div 2^\alpha)={\rm Rem}(k^2\div 2^\alpha),\quad
0\leq k\leq 2^{\alpha-1}.
\end{equation}
$(d)$ For any integers $\ell\geq 0$ and $\alpha\geq 1$ we have
\begin{equation}\label{6.15}
{\rm Rem}((2\ell+1)^2\div 2^{\alpha+2})\equiv 1\pmod{8}.
\end{equation}
\end{lemma}

\begin{lemma}\label{lem:6.2}
$(a)$ For any integer $\beta\geq 1$ we have
\begin{align}
\sum_{k=1}^{2^{2\beta-1}}{\rm Rem}\left(k^2\div 2^{2\beta}\right)
&= 2^{2\beta-2}\left(2^{2\beta}-3\cdot 2^{\beta}+3\right),\label{6.6}\\
\sum_{k=1}^{2^{2\beta}}{\rm Rem}\left(k^2\div 2^{2\beta+1}\right)
&= 2^{2\beta-1}\left(2^{2\beta+1}-4\cdot 2^{\beta}+3\right).\label{6.7}
\end{align}
$(b)$ If $p\equiv 1\pmod{4}$ is a prime, then
\begin{align}
\sum_{k=1}^{\frac{p^{2\beta}-1}{2}}{\rm Rem}\left(k^2\div p^{2\beta}\right)
&=\frac{1}{4}p^{3\beta}\left(p^{\beta}-1\right)\qquad(\beta\geq 1),\label{6.8}\\
\sum_{k=1}^{\frac{p^{2\beta+1}-1}{2}}{\rm Rem}\left(k^2\div p^{2\beta+1}\right)
&=\frac{1}{4}p^{3\beta+1}\left(p^{\beta+1}-1\right)\qquad(\beta\geq 0).\label{6.9}
\end{align}
$(c)$ If $p\equiv 3\pmod{4}$ is a prime, then for $\beta\geq 1$, resp.\
$\beta\geq 0$,
\begin{align}
\sum_{k=1}^{\frac{p^{2\beta}-1}{2}}{\rm Rem}\left(k^2\div p^{2\beta}\right)
&=\frac{1}{4}p^{2\beta}\left(p^{\beta}-1\right)
\left(p^{\beta}-\frac{2}{p-1}h^*(-p)\right),\label{6.10}\\
\sum_{k=1}^{\frac{p^{2\beta+1}-1}{2}}{\rm Rem}\left(k^2\div p^{2\beta+1}\right)
&=\frac{1}{4}p^{2\beta+1}\left(p^{\beta+1}-1\right)
\left(p^{\beta}-\frac{2}{p-1}h^*(-p)\right).\label{6.11}
\end{align}
\end{lemma}

\begin{proof}[Proof of Lemma~\ref{lem:6.2}]
(a) We prove \eqref{6.6} by induction on $\beta$. When $\beta=1$, then 
$\sum_{k=1}^2{\rm Rem}(k^2\div 4)=1+0$, while on the right of \eqref{6.6} we
have $2^0(2^2-3\cdot 2^1+3)=1$; this is the induction beginning. 

For the induction step, we use symmetry (Lemma~\ref{lem:6.3}(c)) to obtain
\begin{align}
&\sum_{k=1}^{2^{2\beta+1}}{\rm Rem}(k^2\div 2^{2\beta+2})
= 2\sum_{k=1}^{2^{2\beta}}{\rm Rem}(k^2\div 2^{2\beta+2})\label{6.16}\\
&\qquad=2\sum_{\ell=1}^{2^{2\beta-1}}{\rm Rem}((2\ell)^2\div 2^{2\beta+2})
+ 2\sum_{\ell=0}^{2^{2\beta-1}-1}{\rm Rem}((2\ell+1)^2\div 2^{2\beta+2}).\nonumber
\end{align}
We deal with the sums in the second row of \eqref{6.16} separately. First, by
\eqref{6.12} we have
\begin{equation}\label{6.17}
\sum_{\ell=1}^{2^{2\beta-1}}{\rm Rem}((2\ell)^2\div 2^{2\beta+2})
= 4\sum_{\ell=1}^{2^{2\beta-1}}{\rm Rem}(\ell^2\div 2^{2\beta})
= 2^{2\beta}\left(2^{2\beta}-3\cdot 2^{\beta}+3\right),
\end{equation}
where we have used \eqref{6.6} as induction hypothesis.

To deal with the second sum on the right of \eqref{6.16}, we note that by parts 
(b) and (d) of Lemma~\ref{lem:6.3} the $2^{2\beta-1}$ summands are distinct 
positive integers of the form $8j+1$, $0\leq j\leq 2^{2\beta-1}-1$. Since
$8j+1 < 8\cdot 2^{2\beta-1} = 2^{2\beta+2}$, these summands are not reduced
modulo $2^{2\beta+2}$, and we have
\begin{align}
\sum_{\ell=0}^{2^{2\beta-1}-1}{\rm Rem}((2\ell+1)^2\div 2^{2\beta+2})
&= \sum_{\ell=0}^{2^{2\beta-1}-1}(8j+1) \label{6.18}\\
&= 2^{2\beta-1}+8\cdot\frac{1}{2}\left(2^{2\beta-1}-1\right)2^{2\beta-1}\nonumber \\
&= 2^{2\beta-1}\left(2^{2\beta+1}-3\right).\nonumber
\end{align}
Finally, substituting \eqref{6.17} and \eqref{6.18} into \eqref{6.16}, we get 
after some straightforward manipulations,
\[
\sum_{k=1}^{2^{2\beta+1}}{\rm Rem}(k^2\div 2^{2\beta+2})
= 2^{2\beta}\left(2^{2\beta+2}-3\cdot 2^{\beta+1}+3\right).
\]
This completes the proof of \eqref{6.6} by induction. The proof of \eqref{6.7}
is analogous; we leave the details to the reader.

(b) We begin with \eqref{6.9}, using again induction on $\beta$. For $\beta=0$,
\eqref{6.9} reduces to 
\begin{equation}\label{6.18a}
\sum_{k=1}^{\frac{p-1}{2}}{\rm Rem}(k^2\div p)=\frac{1}{4}p(p-1),
\end{equation}
which is true by \eqref{2.24} and \eqref{2.25}. We now assume that \eqref{6.9}
holds for some $\beta\geq 0$. Using the fact that for any integer $\alpha\geq 1$
we have the symmetry relation
\[
\left(p^{\alpha}-k\right)^2 = p^{2\alpha}-2kp^{\alpha}+k^2
\equiv k^2\pmod{p^{\alpha}},
\]
we can write
\begin{align}
2\sum_{k=1}^{\frac{p^{2\beta+3}-1}{2}}{\rm Rem}\left(k^2\div p^{2\beta+3}\right)
&=\sum_{k=1}^{p^{2\beta+3}-1}{\rm Rem}\left(k^2\div p^{2\beta+3}\right)\label{6.19}\\
&=\sum_{r=0}^{p-1}\sum_{j=0}^{p^{2\beta+2}-1}
{\rm Rem}\left((jp+r)^2\div p^{2\beta+3}\right).\nonumber
\end{align}
We will now evaluate the inner sum on the right of \eqref{6.19}. When $r=0$,
then \eqref{6.12} gives
\begin{equation}\label{6.20}
\sum_{j=0}^{p^{2\beta+2}-1}{\rm Rem}\left((jp)^2\div p^{2\beta+3}\right)
=p^2\sum_{j=0}^{p^{2\beta+2}-1}{\rm Rem}\left(j^2\div p^{2\beta+1}\right).
\end{equation}
We split $j$ as $j=\ell\cdot p^{2\beta+1}+s$, with $0\leq\ell\leq p-1$ and
$0\leq s\leq p^{2\beta+1}-1$ and note that
\[
{\rm Rem}\left((\ell\cdot p^{2\beta+1}+s)^2\div p^{2\beta+1}\right)
={\rm Rem}\left(s^2\div p^{2\beta+1}\right).
\]
Then
\begin{align*}
\sum_{j=0}^{p^{2\beta+2}-1}{\rm Rem}\left(j^2\div p^{2\beta+1}\right)
&=\sum_{\ell=0}^{p-1}\sum_{s=0}^{p^{2\beta+1}-1}
{\rm Rem}\left(s^2\div p^{2\beta+1}\right) \\
&=2p\sum_{s=0}^{\frac{p^{2\beta+1}-1}{2}}
{\rm Rem}\left(s^2\div p^{2\beta+1}\right) \\
&= 2p\cdot\frac{1}{4}p^{3\beta+1}\left(p^{\beta+1}-1\right),
\end{align*}
where we have used symmetry and then \eqref{6.9} as induction hypothesis.
With \eqref{6.20} we now get
\begin{equation}\label{6.21}
\sum_{j=0}^{p^{2\beta+2}-1}{\rm Rem}\left((jp)^2\div p^{2\beta+3}\right)
= \frac{1}{2}p^{3\beta+4}\left(p^{\beta+1}-1\right).
\end{equation}
Next we consider $1\leq r\leq p-1$ in the inner sum in \eqref{6.19}; for 
greater generality, we consider integers $\alpha\geq 0$. To simplify notation,
we set
\begin{equation}\label{6.22}
a_r := {\rm Rem}(r^2\div p),
\end{equation}
and note that $(jp+r)^2\equiv a_r\pmod{p}$. By Lemma~\ref{lem:6.3}(b), the
terms ${\rm Rem}((jp+r)^2\div p^{\alpha+2})$ are all distinct for
$0\leq j\leq p^{\alpha+1}-1$, they are all congruent to $a_r\pmod{p}$, and lie
between 0 and $p^{\alpha+2}$. So they are $a_r+jp$, 
$0\leq j\leq p^{\alpha+1}-1$, in some order. Hence, for $1\leq r\leq p-1$,
\begin{align*}
S_r&:=\sum_{j=0}^{p^{\alpha+1}-1}{\rm Rem}((jp+r)^2\div p^{\alpha+2})
=\sum_{j=0}^{p^{\alpha+1}-1}(a_r+jp) \\
&= p^{\alpha+1}a_r + \frac{p}{2}\left(p^{\alpha+1}-1\right)p^{\alpha+1}
= p^{\alpha+1}\left(a_r+ \frac{p}{2}\left(p^{\alpha+1}-1\right)\right).
\end{align*}
Next we use the fact that by \eqref{6.22} and \eqref{6.18a} we have
$\sum_{r=1}^{p-1}a_r=\frac{1}{2}p(p-1)$, so that for $\alpha\geq 0$,
\begin{equation}\label{6.23}
\sum_{r=1}^{p-1}S_r
=p^{\alpha+1}\left(\frac{p(p-1)}{2}
+(p-1)\cdot\frac{p}{2}\left(p^{\alpha+1}-1\right)\right)
=\frac{p-1}{2}\cdot p^{\alpha+3}.
\end{equation}
Finally, upon setting $\alpha=2\beta+1$ in \eqref{6.23} and substituting this
and \eqref{6.21} into \eqref{6.19}, we get
\[
\sum_{k=1}^{\frac{p^{2\beta+3}-1}{2}}{\rm Rem}\left(k^2\div p^{2\beta+3}\right)
= \frac{1}{4}\cdot p^{3\beta+4}\left(p^{\beta+2}-1\right).
\]
Comparing this with \eqref{6.9}, we see that the proof by induction is complete.

To prove \eqref{6.8}, we first note that \eqref{6.23} with $\alpha=0$ gives
\[
\sum_{r=1}^{p-1}\sum_{j=0}^{p-1}{\rm Rem}((jp+r)^2\div p^2)
= \frac{p-1}{2}\cdot p^3,
\]
while for $r=0$ we have ${\rm Rem}((jp)^2\div p^2)=0$. Hence
\[
\sum_{k=1}^{p^2-1}{\rm Rem}(k^2\div p^2) = \frac{1}{2}\cdot p^3(p-1),
\]
which is the induction beginning for $\beta=1$ if we take symmetry into 
account. The remainder of the proof of \eqref{6.8} is completely analogous to
that of \eqref{6.9}.

(c) The proofs of the identities \eqref{6.10} and \eqref{6.11} are similar to
those of \eqref{6.8} and \eqref{6.9}, and we leave the details to the 
interested reader. However, the case $p=3$ requires some attention. Rather than
dealing with the details of Dirichlet's class number formula \eqref{3.2}, we
verify that \eqref{3.7} holds for $p=3$ if $h(-3)$ is replaced by 
$\frac{1}{3}=h^*(-3)$. Similarly, \eqref{4.11} holds for $p=3$ with $h(-3)$
replaced by $h^*(-3)$.

This last identity is then the induction beginning, with $\beta=0$, in the
proof of \eqref{6.11}. One other difference between the proofs of \eqref{6.11}
and \eqref{6.9} is that for summing the terms $a_r$ we need to use \eqref{4.11}
again, with the appropriate change for $p=3$. The proof of \eqref{6.10} is
again similar.
\end{proof}

\end{document}
